\documentclass[conference]{IEEEtran}
\usepackage{amsmath}
\usepackage{booktabs}
\usepackage{graphicx}

\title{VALKYRIE: Visual and Audio Luminance Enhancement via Knowledge-driven Yield and Real-time Intelligence Engine}

\author{\IEEEauthorblockN{George David Tsitlauri}
\IEEEauthorblockA{
gdtsitlauri@gmail.com\\
gdtsitlauri.dev\\
github.com/gdtsitlauri
}}

\begin{document}
\maketitle

\begin{abstract}
VALKYRIE is a research-oriented media-enhancement framework targeting budget NVIDIA hardware. The repository combines a modular enhancement pipeline, optional native CUDA kernels, a real-time benchmarking harness, and TensorRT integration code. The current committed evidence strongly supports the engineering claim that the lightweight real-time path can exceed 60 FPS on selected scene classes on a GTX 1650. The quality-benchmark layer is useful for repository-internal diagnostics, but the committed snapshot should be interpreted cautiously because it does not yet constitute a fully standardized full-reference super-resolution benchmark.
\end{abstract}

\section{Introduction}
Open media-enhancement research is often split between high-quality offline systems and fast but opaque application-specific pipelines. VALKYRIE aims to bridge that gap by exposing enhancement modules, benchmarking code, and deployment-oriented optimization hooks in a single repository.

\section{Repository Scope}
The current release includes:
\begin{itemize}
\item a Python package for video and audio enhancement,
\item optional CUDA-native acceleration,
\item benchmarking and reporting utilities,
\item TensorRT conversion and speed-comparison support,
\item committed result summaries for throughput and media benchmarking.
\end{itemize}

\section{Real-Time Evidence}
The strongest current evidence in the repository is the GTX 1650 throughput profile. The committed benchmark summary reports mean frame rates of 78.25 FPS for action scenes and 62.84 FPS for film scenes, with more demanding low-light and game scenarios reducing throughput substantially. These results justify describing VALKYRIE as a credible real-time enhancement engineering project on constrained hardware.

\section{Quality-Benchmark Interpretation}
The repository also contains a media-quality benchmark summary with metrics such as PSNR, SSIM, and LPIPS. However, the committed snapshot is best treated as a repository-internal diagnostic rather than a definitive public leaderboard. The result files document a proxy-reference workflow, and therefore these numbers should be used to track regressions or relative internal changes, not to claim benchmark superiority over canonical super-resolution baselines.

\section{Optimization Path}
TensorRT integration is implemented directly in \texttt{src/valkyrie/tensorrt\_optimize.py}. This is meaningful evidence of deployment-aware optimization work. At the same time, the current repository should still be presented as a prototype optimization path rather than a mature production deployment stack.

\section{Limitations}
Three limitations are important. First, the quality-reporting protocol needs a stricter canonical reference design. Second, the evaluation focuses on local GTX 1650 experiments rather than broader hardware sweeps. Third, no live application deployment traces are included in the current artifact set.

\section{Conclusion}
VALKYRIE is a credible enhancement-engineering repository with real performance evidence, optional TensorRT optimization support, and substantial implementation breadth. Its strongest current contribution is the open real-time pipeline and measurement tooling; stronger full-reference quality claims should wait for a regenerated standardized benchmark suite.

\bibliographystyle{IEEEtran}
\bibliography{references}
\end{document}
